\scanpage{159}\mspage{14}
Cela correspond au phénomène bien connu dans la théorie du corps de \unsure{classe|classes} que la structure du $\pi_{1\,\mathrm{ab}}$ de $X$, quand $X$ est non complète, est \uline{\unsure{continue|infinie}} -- donc il en est de $\bar{\mathrm m}$ du $\mathrm{H}^{1}(X,\mathbb{Z}/p^{n}\mathbb{Z})$ disons.

À ce propos, je te signale que Serre avait proposé (sans jamais l'écrire, je crois) une théorie de dualité des groupes algébriques commutatifs unipotents\inserted{ \unsure{connexes|connexe}}, \deleted{\unsure{duals|doubles|dualité}} pris \uline{mod isogénie radicielle}, dualité à valeurs dans $\mathbb{Q}/\mathbb{Z}$ (ou $\mathbb{Q}_p/\mathbb{Z}_p$). Il trouve que si (sur $k$ alg clos, disons) $G$ est un tel groupe, alors $G'=\operatorname{Ext}^{1}(G,\mathbb{Q}/\mathbb{Z})$ peut être muni \inserted{canoniquement} d'une structure de groupe quasi-algébrique (i.e. défini mod isog. radicielle), de façon unique sans doute en vérifiant qques propriétés de fonctorialité, et en exigeant que pour $G=\mathbb{G}_a$ on trouve sur $\operatorname{Ext}^{1}(\mathbb{G}_a,\mathbb{Q}/\mathbb{Z})\simeq\mathbb{G}_a$ la structure habituelle. Soit $\Delta G=G'=\underline{\mathrm{Ext}}^{1}(G,\mathbb{Q}/\mathbb{Z})$, On trouve $G\simeq\Delta\Delta G$ i.e. $\Delta$ est une authentique autodualité ! J'appelle $\Delta$ la \uline{dualité de Serre}. Elle passe bien sûr aux ind-pro-groupes sur un corps de base quelconque (pas nec alg. clos) \inserted{de car $p>0$.} D'ailleurs, pour des groupes finis étale, c'est $\underline{\mathrm{Ext}}^{0}(G,\mathbb{Q}/\mathbb{Z})$ (dualité de \unsure{Pontrjagin|Pontryagin}) qui fournit une dualité parfaite. On peut \unsure{visser|viser} ensemble, dans une catégorie dérivée convenable, dualité de Serre et dualité de \unsure{Pontrjagin|Pontryagin}, en prenant
% archive sheet number: 159
